\section{Relaciones exactas: propiedades de herencia}
\label{sec:inheritance}

Este método extiende los criterios de equivalencia ya conocidos, la permutación de estados y las reglas espejo, a reglas de espacios diferentes. Por ejemplo, la regla 0 siempre tiene el mismo comportamiento sin importar el número de estados, el tamaño de la vecindad ni la dimensión. Para formalizarlo, llamamos dominio de una regla a la 3-tupla de dimensión, estados y vecindad en la que está definida; dos reglas pertenecen al mismo dominio si comparten los tres elementos.

\subsection{Reglas del mismo dominio}
\label{subsection:rules_same_domain}

En ECA, la permutación de estados y las reglas espejo reducen las 256 reglas a 88 clases~\cite{genaro2013_class}. Sean $\mathcal{A}_1$ y $\mathcal{A}_2$ dos autómatas celulares del mismo dominio y $(x_1, \dots, x_m) \in \mathcal{S}^m$ el estado de la vecindad. Decimos que son equivalentes si se cumple alguna de las siguientes condiciones:
\begin{itemize}
    \item \textbf{Permutación de estados}: existe una biyección $\pi : \mathcal{S} \to \mathcal{S}$ tal que
    \[
    \mathcal{R}_2(x_1, \dots, x_m) = \pi\left( \mathcal{R}_1\left( \pi^{-1}(x_1), \dots, \pi^{-1}(x_m) \right) \right), \quad \forall (x_1, \dots, x_m) \in \mathcal{S}^m,
    \]
    es decir, si al renombrar los estados la estructura sigue siendo la misma. Para $\mathcal{S} = \{0, 1\}$, la única permutación no trivial es $\pi(x) = 1 - x$.

    \item \textbf{Reglas espejo}: sea $M(x_1, \dots, x_m) = (x_m, \dots, x_1)$ la reflexión de la vecindad; $\mathcal{A}_2$ es la regla espejo de $\mathcal{A}_1$ si $\mathcal{R}_2 = \mathcal{R}_1 \circ M$. Es decir, si reflejamos la configuración inicial, la evolucionamos con $\mathcal{A}_1$ y reflejamos el resultado, obtenemos la evolución de $\mathcal{A}_2$.
\end{itemize}

\subsection{Reducción de estados}

Sea $\mathcal{A}$ un autómata celular y $a \in \mathcal{S}$ uno de sus estados. La reducción de estados es posible cuando:
\begin{enumerate}
    \item El estado $a$ no es producido por la regla: $\mathcal{R}(x_1, \dots, x_m) \neq a, \; \forall (x_1, \dots, x_m) \in \mathcal{S}^m$.
    \item Toda vecindad que contiene al estado $a$ produce siempre el mismo estado $b$: $\mathcal{R}(x_1, \dots, x_m) = b$ si $x_j = a$ para algún $j$.
\end{enumerate}
Entonces el autómata reducido es $\mathcal{A}' = (\mathcal{L}, \mathcal{S}', \mathcal{N}, \mathcal{R}')$ con $\mathcal{S}' = \mathcal{S} \setminus \{a\}$ y $\mathcal{R}'$ la restricción de $\mathcal{R}$ a $\mathcal{S}'^m$. Como el estado $a$ nunca es producido, solo puede aparecer en la condición inicial, y como toda vecindad que lo contiene produce $b$, desaparece en el primer paso; a partir de ahí, la evolución de $\mathcal{A}$ es exactamente la de $\mathcal{A}'$. En sentido inverso, toda regla de $q$ estados se puede extender a $q + 1$ estados agregando un estado nuevo que nunca se produce y que envía al 0 toda vecindad en la que aparece.

\subsection{Reducción de vecindad}

La reducción de vecindad es posible cuando el estado de un vecino no afecta a la transición, es decir, existe un índice $j$ tal que
\[
\mathcal{R}(x_1, \dots, x_j, \dots, x_m) = \mathcal{R}(x_1, \dots, x_j', \dots, x_m), \quad \forall (x_1, \dots, x_m) \in \mathcal{S}^m, \; \forall x_j' \in \mathcal{S}.
\]
Entonces el autómata reducido es $\mathcal{A}' = (\mathcal{L}, \mathcal{S}, \mathcal{N}', \mathcal{R}')$ con $\mathcal{N}' = \mathcal{N} \setminus \{\vec{v}_j\}$. Si el vecino eliminado no está en uno de los extremos, la vecindad se ajusta para que sea contigua de $m - 1$ células, bajo la hipótesis de que la distancia entre los vecinos solo transforma el espacio de evolución y no el comportamiento.

Aplicando estas cuatro relaciones de forma iterativa se construye un grafo de relaciones entre reglas de dominios distintos, cuyas componentes conexas son las clases de equivalencia; en ECA quedan 81 clases. Aunque esta clasificación sigue siendo binaria, funciona como referencia para validar los métodos siguientes.
